\subsubsection{The corrected global spectrum}
\label{sec:p10-corrected-global-spectrum}

The construction uses a site presented by finite lattice relations,
not the object of already existing prime predicates.
Throughout, a local ring means the internal notion of
\cref{def:p10-lawvere-sheaves-prime}. A homomorphism between local rings
is called local when it reflects units.
All finite constructions below are internal to the base topos.
They require no natural numbers object.

\begin{hbdefinition}{def:p10-spectrum-finite-subobjects}{Finite subobjects for lattice constructions}
For an object \(X\), let \(K(X)\) be the least subobject of its
power object containing the empty subobject and every singleton,
and closed under binary union. Leastness is defined by intersecting
all subobjects of the power object with these properties.
Members of \(K(X)\) are the Kuratowski-finite subobjects used here.
They need not have decidable equality or a chosen enumeration.
\end{hbdefinition}

\begin{hbtheorem}{thm:p10-spectrum-lattice}{proved=annex:p10-spectrum-lattice}{The finite lattice presentation}
There is a bounded distributive lattice \(L_A\) with generators
\(d(a)\), internally indexed by \(a\in A\), and relations
\[
 d(0)=0,\qquad d(1)=1,\qquad
 d(ab)=d(a)\wedge d(b),\qquad
 d(a+b)\leq d(a)\vee d(b).
\]
It is universal for maps from \(A\) into bounded distributive lattices
satisfying these relations.
For every geometric morphism \(f:\mathcal F\to\mathcal E\),
there is a canonical isomorphism
\(f^*L_A\cong L_{f^*A}\) preserving the generators.
The frame of ideals of \(L_A\) is canonically isomorphic to the
radical-ideal frame of \cref{thm:p10-radical-spectrum-frame}.
\end{hbtheorem}

Regard \(L_A\) as a thin category, with one arrow \(v\to u\)
when \(v\leq u\). A basic cover of \(u\) is a
Kuratowski-finite subobject \(V\subseteq{\downarrow}u\) with
\(\bigvee V=u\). A sieve covers when it contains a basic cover.
The empty family covers the bottom element.

\begin{hbtheorem}{thm:p10-spectrum-site-classification}{proved=annex:p10-spectrum-site-classification}{Classification by the finite-cover site}
These covers define an internal Grothendieck topology.
Its sheaf category \(\mathcal S_A\) is a topos with geometric
projection \(p:\mathcal S_A\to\mathcal E\).
For any geometric morphism \(f:\mathcal F\to\mathcal E\),
geometric morphisms \(g:\mathcal F\to\mathcal S_A\) over
\(\mathcal E\), up to compatible isomorphism, correspond to
bounded distributive lattice homomorphisms
\[
 f^*L_A\longrightarrow\Omega_{\mathcal F},
\]
equivalently to prime predicates on \(f^*A\).
Here \(\Omega_{\mathcal F}\) denotes the target's truth-value frame.
The corresponding prime predicate is the inverse image of a
generic subobject \(S\hookrightarrow p^*A\) in \(\mathcal S_A\).
\end{hbtheorem}

\begin{hbtheorem}{thm:p10-corrected-global-spectrum}{proved=annex:p10-corrected-global-spectrum}{Universal local ring}
Put \(\mathcal O=S^{-1}(p^*A)\), where \(S\) is the generic
prime predicate. Then \(\mathcal O\) is a local ring object.
For every geometric morphism \(f:\mathcal F\to\mathcal E\),
local ring object \(B\) of \(\mathcal F\), and unital ring
homomorphism \(h:f^*A\to B\), there is a factorization
\[
 f=p\,g,\qquad
 f^*A\longrightarrow g^*\mathcal O
          \xrightarrow{\ k\ }B
\]
whose second map is local and whose composite is \(h\).
The factorization is unique up to the unique compatible
isomorphism over \(\mathcal E\).
This is the corrected global spectrum of the ringed topos
\((\mathcal E,A)\).
\end{hbtheorem}

The distinction from the original construction is structural.
The site records the finite relations that a prime predicate must
satisfy; it does not replace those relations by their truth at primes
already present in \(\mathcal E\). Its generic predicate exists in
the new sheaf topos. The preceding evaluation-map counterexample
therefore does not collapse this construction.
